\section{Literature Review}

\subsection{Academic Performance}

Academic performance is commonly used as an indicator of students' educational outcomes. It may be measured using examination scores, continuous assessment results, grade point averages, or other academic indicators \citep{determinants2023academic}.

\subsection{Factors Associated with Academic Performance}

Several factors may influence students' academic outcomes. These factors can be broadly classified into academic, personal, social, and demographic characteristics \citep{determinants2023academic}. Study habits, study skills, and attitudes towards learning have also been identified as important predictors of academic performance \citep{crede2008study}.

\subsection{Study Time and Performance}

Study time represents the amount of time a student spends engaging with academic materials outside scheduled classes. Although study time is an important aspect of students' academic behaviour, research suggests that the amount of time spent studying does not necessarily translate directly into higher academic performance. The effectiveness and quality of study practices may be important considerations alongside the amount of time spent studying \citep{plant2005study}.

\subsection{Attendance and Performance}

Class attendance is another measurable factor that may be associated with academic performance. Research has examined the relationship between attendance and academic outcomes, with findings suggesting that attendance can be an important predictor of student performance \citep{arnold2018attendance}. The timing and pattern of student absences may also influence the relationship between absenteeism and academic achievement \citep{keppens2023absence}.

\subsection{Statistical Modelling}

Statistical modelling provides a framework for examining relationships between a dependent variable and one or more explanatory variables. Multiple linear regression can be used to model academic performance as a function of measurable student characteristics. A general multiple linear regression model can be expressed as

\[
Y = \beta_0 + \beta_1X_1 + \beta_2X_2 + \cdots + \beta_pX_p + \varepsilon,
\]

where \(Y\) represents academic performance, \(X_1,\ldots,X_p\) represent explanatory variables, \(\beta_0\) is the intercept, \(\beta_1,\ldots,\beta_p\) are regression coefficients, and \(\varepsilon\) is the error term. Previous research has demonstrated the use of multiple linear regression and principal component analysis for predicting students' academic performance \citep{yang2018predicting}.